% ============================================================================
% 第 4 章  分裂生成与扭转  （原书 pp.51–60，PDF pp.54–63）
% 底本：Alessandro Imparato, Towards Homological Mirror Symmetry, ETH Zürich
% 说明：原书公式全书连续编号，本章恰为 (52)–(53) 共 2 个编号公式；main.tex 采用
%       全局连续编号，故自动编号与原文一致。每个编号公式恰用一个 equation 环境
%       （内含 tikzcd，只编号一次），其余公式一律 equation*/\[ \] 不编号；\label
%       按原书编号命名（eq:4.52、eq:4.53）。
% ============================================================================
\section{分裂生成与扭转}\label{ch:04}

\takeaway{%
  核心图像：幂等态射的“像”在普通范畴里给出分裂封闭与 Karoubi 完备化，在
  $A_\infty$ 世界里则对应同伦幂等与抽象像；补进一切幂等分裂，就得到分裂封闭导出范畴
  $\mathsf D^\pi(\mathcal A)$，它是 Fukaya 范畴“分裂生成”讨论的舞台。\\
  数学工具：幂等自同态与收缩—包含对、Karoubi 完备化、同伦幂等 $\wp$、抽象像模
  $\hom_{\mathcal A}(X,Y)[q]$、三角包络 $\mathsf{Tw}\mathcal A$、抽象扭转函子
  $\mathcal T_Y$ 及其伴随 $\mathcal T^\vee_{Y_0}$。\\
  主要结论：每个 $A_\infty$-范畴都有分裂封闭（命题 4.13）；三角 $A_\infty$-范畴的分裂
  封闭仍三角（引理 4.14）；扭转可迭代复合，其零性或平移性构成分裂生成的判据
  （推论 4.24）。%
}

\subsection{范畴论预备}\label{sec:4.1}

\begin{Def}\label{def:4.1}
设 $\mathsf C$ 为 $\mathbb K$ 上的线性范畴。考虑某个 $Y\in\operatorname{obj}(\mathsf C)$
的\emph{幂等自同态}（idempotent endomorphism）$p$，即满足 $p\circ p=p$ 的态射
$p\in\Hom_{\mathsf C}(Y,Y)$。

若三元组 $(Z,r,k)$（其中 $Z\in\operatorname{obj}(\mathsf C)$，
$r\in\Hom_{\mathsf C}(Y,Z)$ 称为\emph{收缩}（retraction），
$k\in\Hom_{\mathsf C}(Z,Y)$ 称为\emph{包含}（inclusion））满足
\[
  r\circ k=\mathrm{id}_Z
  \qquad\text{且}\qquad
  k\circ r=p,
\]
则称之为 $p$ 的\emph{像}（image）。相应地，若这样的三元组存在，则称 $p$
\emph{分裂}（splits）。我们把上述数据编码为下图：
\begin{equation}\label{eq:4.52}
  \begin{tikzcd}
    Z \arrow[loop left, "{\mathrm{id}_Z}"] \arrow[r, bend left=25, "k"]
      & Y \arrow[l, bend left=25, "r"'] \arrow[loop right, "p"]
  \end{tikzcd}
\end{equation}
若 $\mathsf C$ 的所有幂等自同态都分裂，则称 $\mathsf C$ 为\emph{分裂封闭的}
（split-closed）。
\end{Def}

注意，幂等态射 $p$ 的像若存在则唯一：给定另一个合适的三元组 $(Z',r',k')$，
$r'\circ k:Z\to Z'$ 是对象间的同构，其逆为 $r\circ k'$。
% [Tutor Note: 【易错】这里的“像”不是核—余核意义下的像，而是把 $p$ 显式分解为“包含—收缩”
% 的数据 $(Z,r,k)$；三条关系缺一不可，且方向极易写反：$r:Y\to Z$ 是收缩（先作用），
% $k:Z\to Y$ 是包含。图中 $k$ 在上、$r$ 在下，与式 (52) 的编号约定一致。]

还要注意，若 $\mathsf C$ 是加性范畴（例如三角范畴）且分裂封闭，则对幂等态射
$p\in\Hom_{\mathsf C}(Y,Y)$，可以取对 $Z:=\operatorname{im}(p)$ 与
$Z^{\perp}:=\operatorname{im}(\mathrm{id}_Y-p)$，得到分解 $Y=Z\oplus Z^{\perp}$。

\begin{Def}\label{def:4.2}
范畴 $\mathsf C$ 的\emph{分裂封闭}（split-closure）是一对 $(\mathsf D,F)$：$\mathsf D$
为分裂封闭范畴，$F:\mathsf C\to\mathsf D$ 为全忠实嵌入，且 $\mathsf D$ 的任一对象
$Z$ 都是某个幂等自同态 $p\in\Hom_{\mathsf D}(F(X),F(X))$ 的像，其中
$X\in\operatorname{obj}(\mathsf C)$。
\end{Def}

\begin{Rem}\label{rem:4.3}
任何范畴 $\mathsf C$ 都有一个自动给出的分裂封闭 $\overline{\mathsf C}$，称为
\emph{Karoubi 完备化}（Karoubi completion），也称\emph{幂等完备化}（idempotent
completion）。其定义如下：对象为二元组 $(Y,p)$，其中 $Y\in\operatorname{obj}(\mathsf C)$
而 $p\in\Hom_{\mathsf C}(Y,Y)$ 为幂等自同态（$\mathsf C$ 的每个对象都至少有一个这样的
幂等自同态，即 $\mathrm{id}_Y$）；态射为
\[
  \Hom_{\overline{\mathsf C}}\bigl((Y_0,p_0),(Y_1,p_1)\bigr)
  :=p_1\circ\Hom_{\mathsf C}(Y_0,Y_1)\circ p_0,
\]
即满足 $f=p_1\circ f\circ p_0$ 的三元组 $(p_0,f,p_1):=p_1\circ f\circ p_0$，其中
$f\in\Hom_{\mathsf C}(Y_0,Y_1)$。复合是由 $\mathsf C$ 诱导的显然复合，恒等态射均取
\[
  \mathrm{id}_{(Y,p)}:=(p,p,p)=p.
\]
全忠实嵌入为 $E:\mathsf C\to\overline{\mathsf C}$，
$Y\mapsto(Y,\mathrm{id}_Y)$，$p\mapsto(\mathrm{id}_Y,p,\mathrm{id}_Y)$。

确实，$\overline{\mathsf C}$ 分裂封闭：对任一幂等态射
$\bar p=(p,s,p)\in\Hom_{\overline{\mathsf C}}\bigl((Y,p),(Y,p)\bigr)$（其中
$s\in\Hom_{\mathsf C}(Y,Y)$ 必然满足 $s=p\circ s\circ p$ 与 $s=s\circ p\circ s$），取
$Z:=(Y,\mathrm{id}_Y)$、$r:=(p,s,\mathrm{id}_Y):(Y,p)\to Z$ 与
$k:=(\mathrm{id}_Y,s,p):Z\to(Y,p)$，便有
\[
  r\circ k=\mathrm{id}_Y\circ s\circ\mathrm{id}_Y=E(s),
  \qquad
  k\circ r=p\circ s\circ p=\bar p.
\]
事实上可以证明，$\mathsf C$ 的任意分裂封闭 $(\mathsf D,F)$ 都等价于 Karoubi 完备化
$(\overline{\mathsf C},E)$。此外，若 $\mathsf C$ 是三角范畴，则 $\overline{\mathsf C}$
连同嵌入 $E$ 也是三角范畴。下面给出一个更具体的结果。
\end{Rem}
% [Tutor Note: 【背景】Karoubi 完备化的对象应理解为“形式上补上 $p$ 的像”；把 $(Y,p)$ 看作
% $\operatorname{im}(p)$，态射 $p_1\circ f\circ p_0$ 就是“先在 $\operatorname{im}(p_0)$ 上取值、
% 再落到 $\operatorname{im}(p_1)$”的诱导映射。这是幂等完备化的标准构造；可参看
% \cite{GM03} 第 II 章关于加性范畴的讨论。]

\begin{Lem}\label{lem:4.4}
设 $\mathsf B=(\mathsf B,T)$ 为分裂封闭的三角范畴，$\mathsf A\subset\mathsf B$ 为三角
全子范畴（从而对平移与正合三角形封闭）。设 $\tilde{\mathsf A}\subset\mathsf B$ 为所有
同构于 $\mathsf A$ 中某对象的直和项的对象构成的全子范畴（特别地，
$\tilde{\mathsf A}\cong\overline{\mathsf A}$，因为后者分裂封闭）。则 $\tilde{\mathsf A}$
是三角范畴。
\end{Lem}

\begin{proof}
首先注意，我们处理的是加性范畴，因此谈论直和分解才有意义（参看上文）。按假设，任一
$\tilde Y_i\in\operatorname{obj}(\tilde{\mathsf A})$ 都出现在某个
$Y\cong\tilde Y_i\oplus\tilde Y_i^{\perp}\in\operatorname{obj}(\mathsf A)$ 中。于是可以
取对应的收缩与包含 $r_i,k_i$，使得 $r_i\circ k_j=\mathrm{id}_{\tilde Y_i}$。
% [Tutor Note: 【勘误】原文此处印作 $r_i\circ k_j=\mathrm{id}_{\tilde Y_i}$，下标 $j$
% 显然应为 $i$：只有 $r_i$ 与 $k_i$ 才是一对相容的收缩—包含，$r_i\circ k_i=\mathrm{id}_{\tilde Y_i}$；
% 若真按 $j\ne i$ 理解，该等式一般不成立。译文照印原式，不改动。]

由三角范畴的公理 (T1)（见定义 \ref{def:3.2}），可以把态射
$\tilde c:\tilde Y_0\to\tilde Y_1$ 与
$c:=k_1\circ\tilde c\circ r_0:Y_0\to Y_1$ 分别补全为正合三角形，其中涉及某个第三对象
$\tilde Y_2$ 与 $Y_2$。它们构成如引理 \ref{lem:3.3} 中的两个图（一个以包含
$k_0,k_1$ 为竖直箭头，另一个以收缩 $r_0,r_1$ 为竖直箭头），其最左边的方块由 $c$ 的
定义而交换。因此可以找到补充态射 $k_2:\tilde Y_2\to Y_2$ 与
$r_2:Y_2\to\tilde Y_2$，使整个图交换。把它们复合，得到：
\[
  \begin{tikzcd}[row sep=large]
    \tilde Y_0 \arrow[r, "\tilde c"] \arrow[d, "r_0\circ k_0"']
      & \tilde Y_1 \arrow[r] \arrow[d, "r_1\circ k_1"']
      & \tilde Y_2 \arrow[r] \arrow[d, dashed, "r_2\circ k_2"']
      & T(\tilde Y_0) \arrow[d, "T(r_0\circ k_0)"'] \\
    \tilde Y_0 \arrow[r, "\tilde c"']
      & \tilde Y_1 \arrow[r]
      & \tilde Y_2 \arrow[r]
      & T(\tilde Y_0)
  \end{tikzcd}
\]
由于前两个竖直箭头是同构（它们都是恒等态射），再次应用引理 \ref{lem:3.3} 可知
$r_2\circ k_2$ 也是同构，从而可逆；此外图的交换性迫使它等于恒等态射
$\mathrm{id}_{\tilde Y_2}$。于是取幂等自同态 $p:=k_2\circ r_2\in\Hom_{\mathsf A}(Y_2,Y_2)$（它显然幂等），便看出
$p$ 经由 $(\tilde Y_2,r_2,k_2)$ 分裂，从而 $\tilde Y_2$
同构于 $Y_2$ 的直和项。按定义，这意味着 $\tilde Y_2\in\operatorname{obj}(\tilde{\mathsf A})$，
故 $\tilde{\mathsf A}$ 完全包含上图所示的正合三角形。因此 $\tilde{\mathsf A}$ 是三角
范畴。
\end{proof}
% [Tutor Note: 【易错】图中第二行的三个对象仍是带波浪号的 $\tilde Y_i$，不是 $Y_i$：
% 把“包含”$k_i$ 与“收缩”$r_i$ 复合后又回到了 $\tilde{\mathsf A}$ 一侧，竖直箭头
% $r_i\circ k_i$ 是 $\tilde Y_i$ 的自同态；正是这一点让第二步引理 3.3 能给出
% $r_2\circ k_2$ 的可逆性。]

\subsection{抽象像}\label{sec:4.2}

下面把刚刚得到的理论翻译成 $A_\infty$-模的语言。为此，把基域 $\mathbb K$ 视为平凡的
$A_\infty$-范畴：$\operatorname{obj}(\mathbb K):=\{*\}$，
$\hom_{\mathbb K}(*,*):=\{\mathrm{id}_*\}$，从而与它的上同调范畴一致，
$H(\mathbb K)\cong\mathbb K$。

\begin{Def}\label{def:4.5}
设 $\mathcal A$ 为 $A_\infty$-范畴，$\mathcal Q=\mathrm{mod}(\mathcal A)$。
$Y\in\operatorname{obj}(\mathcal A)$ 的\emph{同伦幂等}（idempotent up to homotopy）
是指一个非幺 $A_\infty$-函子 $\wp:\mathbb K\to\mathcal A$，满足 $\wp(*)=Y$，且对
$d\ge1$ 有
\[
  \wp^d\equiv\wp^d(\mathrm{id}_*,\dots,\mathrm{id}_*)\in\hom^{1-d}_{\mathcal A}(Y,Y),
\]
并满足相应的方程 \eqref{eq:2.10}：
\[
  \sum_{\substack{1\le r\le d\\ s_1+\dots+s_r=d}}
    \mu^r_{\mathcal A}(\wp^{s_r},\dots,\wp^{s_1})
  =
  \begin{cases}
    \wp^{d-1} & \text{若 $d$ 为偶数},\\
    0 & \text{若 $d$ 为奇数},
  \end{cases}
\]
其中的等式落在 $\hom^{2-d}_{\mathcal A}(Y,Y)$ 中。
\end{Def}
% [Tutor Note: 【易错】$\wp^d$ 的次数为 $1-d$：$\wp^1\in\hom^0$、$\wp^2\in\hom^{-1}$，
% 依次递减。左端 $\mu^r_{\mathcal A}$ 把 $r$ 个次数为 $1-s_i$ 的元素相乘，总次数为
% $r-d$，加上 $\mu^r$ 的 $2-r$ 次恰为 $2-d$，与右端 $\wp^{d-1}$ 的次数一致；
% 这是检验符号约定是否抄对的最快办法。]

可以证明幂等自同态与同伦幂等之间有如下关系（见 \cite{Sei08} 引理 4.2）。

\begin{Lem}\label{lem:4.6}
设 $p\in\Hom^0_{H(\mathcal A)}(Y,Y)$ 幂等，$Y\in\operatorname{obj}(\mathcal A)$。则
存在 $Y$ 的一个同伦幂等 $\wp:\mathbb K\to\mathcal A$，使得 $\langle\wp^1\rangle=p$。
\end{Lem}

我们着手把 $\wp$ 描述为某个 $A_\infty$-模的像，做法与上一节处理幂等态射 $p$ 时类似。

\begin{Def}\label{def:4.7}
设 $\wp:\mathbb K\to\mathcal A$ 为某个 $Y\in\operatorname{obj}(\mathcal A)$ 的同伦幂等。
定义 $\wp$ 的\emph{抽象像}（abstract image）为如下 $A_\infty$-模
$\mathcal Z\in\operatorname{obj}(\mathcal Q)$：它对任一
$X\in\operatorname{obj}(\mathcal A)$ 给出
\[
  \mathcal Z(X):=\hom_{\mathcal A}(X,Y)[q],
\]
即关于一个次数为 $-1$ 的变量 $q$ 的形式多项式的分次向量空间。记 $\delta_q$ 为形式
微分算子，使得 $\delta_q(q^0)=0$、$\delta_q(q^k)=q^{k-1}$，则复合映射为
\[
  \mu^1_{\mathcal Z}\bigl(b(q)\bigr)
  :=\sum_{\substack{r\ge2\\ s_2,\dots,s_r\ge1}}
    \delta_q^{s_2+\dots+s_r}\mu^r_{\mathcal A}
    (\wp^{s_r},\dots,\wp^{s_2},b(q))
  +(-1)^{|b|}\frac{b(q)-b(-q)}{2q}
  \in\mathcal Z(X_0)[1],
\]
\[
  \begin{aligned}
  \mu^d_{\mathcal Z}:{}&\mathcal Z(X_{d-1})\otimes
    \hom_{\mathcal A}(X_{d-2},X_{d-1})\otimes\dots\otimes
    \hom_{\mathcal A}(X_0,X_1)\longrightarrow\mathcal Z(X_0)[2-d],\\
  &\bigl(b(q),a_{d-1},\dots,a_1\bigr)\longmapsto
    \sum_{\substack{r\ge d+1\\ s_{d+1},\dots,s_r\ge1}}
    \delta_q^{s_{d+1}+\dots+s_r}\mu^r_{\mathcal A}
    (\wp^{s_r},\dots,\wp^{s_{d+1}},b(q),a_{d-1},\dots,a_1),
  \end{aligned}
\]
对 $d>1$ 成立。
\end{Def}
% [Tutor Note: 【背景】$\mathcal Z(X)=\hom_{\mathcal A}(X,Y)[q]$ 把形式多项式 $q$ 当作
% “把 $\wp$ 的各阶分量逐次插入复合”的记账工具：$\delta_q$ 每作用一次就消耗掉一个 $q$，
% 从而把 $\mu^r_{\mathcal A}$ 的高阶项按 $s_2+\dots+s_r$ 分次。最后一项
% $\frac{b(q)-b(-q)}{2q}$ 是投影到 $q$ 的偶次部分之外，保证 $\mu^1_{\mathcal Z}$ 在形式上
% 只保留与 $b(q)$ 的“奇偶相容”部分；参见 \cite{Sei08} 第 (4c) 节。]

借助 $\mathcal Z(X)$ 的滤过并研究相应的谱序列，最终可在上同调模层面得到如下刻画
（回忆定义 \ref{def:2.22}）。

\begin{Lem}\label{lem:4.8}
考虑某个 $Y\in\operatorname{obj}(\mathcal A)$ 的同伦幂等 $\wp$ 的抽象像 $\mathcal Z$，
并取 $p\in\Hom^0_{H(\mathcal A)}(Y,Y)$ 使 $\langle\wp^1\rangle=p$ 幂等（由引理
\ref{lem:4.6} 可行）。
\begin{enumerate}[label=\textup{\roman*}.,leftmargin=*]
  \item 上同调模 $H(\mathcal Z)\in\operatorname{obj}\bigl(\mathrm{Mod}^0(H(\mathcal A))\bigr)$%
        \footnote{由于我们隐含假设了 $\mathcal A$ 是 c-幺的，$H(\mathcal Z):
        H(\mathcal A)^{opp}\to\mathrm{GradVect}$ 确实是反变幺函子，即
        $\mathrm{Mod}^0(H(\mathcal A))$ 的对象（参看定义 \ref{def:2.22}）。}
        描述如下：对任一 $X\in\operatorname{obj}(\mathcal A)$，
        \[
          H(\mathcal Z)(X)=H\bigl(\mathcal Z(X)\bigr)
          \cong p\circ\Hom_{H(\mathcal A)}(X,Y)
          \in\operatorname{obj}(\mathrm{GradVect}),
        \]
        特别地，$H(\mathcal Z)$ 是 $H(\Upsilon(Y))$ 的直和项；其模乘法为
        \[
          H(\mathcal Z)(X_1)\otimes\Hom_{H(\mathcal A)}(X_0,X_1)
          \to H(\mathcal Z)(X_0),\qquad
          (p\circ\langle a_1\rangle,\langle a_0\rangle)
          \mapsto p\circ(\langle a_1\rangle\circ\langle a_0\rangle).
        \]
  \item 对任一 $\mathcal M\in\operatorname{obj}(\mathcal Q)$，存在分次群同构
        \[
          \Hom_{H(\mathcal Q)}(\mathcal Z,\mathcal M)\cong H(\mathcal M)(Y)\circ p,
        \]
        它关于与模同态的左侧复合是自然的。结合上一点可知
        $\Hom_{H(\mathcal Q)}(\mathcal Z,\mathcal Z)
        \cong p\circ\Hom_{H(\mathcal A)}(Y,Y)\circ p$。
\end{enumerate}
\end{Lem}
% [Tutor Note: 【易错】第 (i) 点里的 $\circ p$ 是右乘 $p$：$p\circ\Hom(X,Y)$ 表示
% “先由 $p$ 投影再取 $\Hom$”，即 $\Hom_{H(\mathcal A)}(X,Y)$ 中落在 $\operatorname{im}(p)$
% 的那些元素；不要与左乘 $p$ 混为一谈。（ii）中的 $H(\mathcal M)(Y)\circ p$ 同理。
% 这一约定是后文 $r\circ p$、$k$ 互为逆的关键。]

\begin{Cor}\label{cor:4.9}
考虑某个 $Y\in\operatorname{obj}(\mathcal A)$ 的同伦幂等 $\wp$ 的抽象像 $\mathcal Z$，
并取 $p\in\Hom^0_{H(\mathcal A)}(Y,Y)$ 使 $\langle\wp^1\rangle=p$ 幂等。则 $\mathcal Z$
是 $\langle\Upsilon^1(\wp^1)\rangle$ 在 $H^0(\mathcal Q)$ 中的像（按定义 \ref{def:4.1}
的意义）。特别地，$\mathcal Z$ 在同构意义下只依赖于类 $p$，而与它的代表 $\wp^1$ 无关。
\end{Cor}

\begin{proof}
记 $\mathcal Y=\Upsilon(Y)\in\operatorname{obj}(\mathcal Q)$ 为 Yoneda 模，注意
$p\in H(\mathcal Y(Y))=H(\mathcal Y)(Y)$。我们验证类
$\langle\Upsilon^1(\wp^1)\rangle=H(\Upsilon)\langle\wp^1\rangle=H(\Upsilon)(p)
\in\Hom^0_{H(\mathcal Q)}(\mathcal Y,\mathcal Y)$ 是 $\mathcal Y$ 在 $H^0(\mathcal Q)$
中的幂等自同态。

这其实是 $H(\Upsilon)$ 的函子性与 $p$ 的幂等性的直接推论：
$H(\Upsilon)(p)\circ H(\Upsilon)(p)=H(\Upsilon)(p\circ p)=H(\Upsilon)(p)$。把
$H(\Upsilon)$ 作用于图 (52)，得到
\begin{equation}\label{eq:4.53}
  \begin{tikzcd}
    \Upsilon(Z) \arrow[loop left, "{\mathrm{id}_{\Upsilon(Z)}}"]
      \arrow[r, bend left=25, "\kappa"]
      & \mathcal Y \arrow[l, bend left=25, "\rho"']
      \arrow[loop right, "{\langle\Upsilon^1(\wp^1)\rangle}"]
  \end{tikzcd}
\end{equation}
上述图位于 $H^0(\mathcal Q)$ 中，其中我们置 $\rho:=H(\Upsilon)(r)$ 与
$\kappa:=H(\Upsilon)(k)$。回忆 $\mathrm{id}_{\Upsilon(Z)}$ 由严格单位
$e_{\Upsilon(Z)}\in\hom_{\mathcal Q}(\Upsilon(Z),\Upsilon(Z))$ 诱导，如第
\ref{sec:2.8} 节所述。像的两条刻画关系显然成立，这仍归功于 $H(\Upsilon)$ 的函子性。
% [Tutor Note: 【勘误】原文此处引“Section 2.8”；但严格单位 $e_{\Upsilon(Z)}$ 的讨论
% 在本书第 2.7 节“严格幺性与上同调幺性”，第 2.8 节是“拟等价”。疑为误引 2.7 为 2.8。
% 译文照印 2.8。]

还需说明 $\mathcal Z\cong\Upsilon(Z)$ 在 $H^0(\mathcal Q)$ 中成立。由引理 \ref{lem:4.8}，
\[
  \Hom_{H(\mathcal Q)}\bigl(\mathcal Z,\Upsilon(Z)\bigr)
  \cong H\bigl(\Upsilon(Z)(Y)\bigr)\circ p
  =\Hom^0_{H(\mathcal A)}(Y,Z)\circ p.
\]
于是只需右端存在同构：一个这样的同构由 $r\circ p$ 给出，其逆为
$k\in\Hom^0_{H(\mathcal A)}(Z,Y)$。
\end{proof}
% [Tutor Note: 【含糊】此证明把 $H(\Upsilon)$ 作用于图 (52)，等于默认 $p$ 在
% $H^0(\mathcal A)$ 中已经分裂（图中 $Z$ 是 $\mathcal A$ 的对象，$r,k$ 为该分裂的数据），
% 而推论 4.9 的假设只要求 $p$ 幂等；引理 4.11 的证明正是先假设 $p$ 分裂、再引用本条。
% 原文未声明这一附加假设，若要处理一般幂等，需直接验证 $(\mathcal Z,\rho,\kappa)$
% 的两条关系。]

作为推论，$H^0(\mathcal Q)$ 是分裂封闭的普通范畴：$H^0(\mathcal Q)$ 中的任一幂等
自同态都（再次由 $H(\Upsilon)$ 的双射性与函子性）对应到 $H^0(\mathcal A)$ 中的一个幂等
自同态，而后者在 $H^0(\mathcal Q)$ 中的像刚由推论 \ref{cor:4.9} 证明分裂。

\subsection{分裂封闭导出范畴}\label{sec:4.3}

\begin{Def}\label{def:4.10}
称 $A_\infty$-范畴 $\mathcal A$ 为\emph{分裂封闭的}（split-closed），若任意同伦幂等的
抽象像 $\mathcal Z$ 都被某个 $Y\in\operatorname{obj}(\mathcal A)$ 准表示，即可以找到
$H^0(\mathcal Q)$ 中的同构 $\langle t\rangle:\Upsilon(Y)=\mathcal Y\xrightarrow{\ \sim\ }
\mathcal Z$。

$\mathcal A$ 的\emph{分裂封闭}（split-closure）是一对 $(\mathcal B,\mathcal F)$，其中
$\mathcal B$ 为分裂封闭 $A_\infty$-范畴，$\mathcal F:\mathcal A\to\mathcal B$ 为上同调
全忠实的 $A_\infty$-函子，且对任一 $Z\in\operatorname{obj}(\mathcal B)$，存在
$Y\in\operatorname{obj}(\mathcal A)$，使 $Z$ 成为某个幂等自同态
$p\in\Hom_{H(\mathcal B)}(\mathcal F(Y),\mathcal F(Y))$ 在 $H^0(\mathcal B)$ 中的像
（经典意义下）。
\end{Def}

\begin{Lem}\label{lem:4.11}
$A_\infty$-范畴 $\mathcal A$ 分裂封闭，当且仅当 $H^0(\mathcal A)$ 分裂封闭（经典意义下）。
\end{Lem}

\begin{proof}
若 $H^0(\mathcal A)$ 分裂封闭，给定同伦幂等 $\wp:\mathbb K\to\mathcal A$，
$\wp(*)=Y\in\operatorname{obj}(\mathcal A)$，先由引理 \ref{lem:4.6} 找到幂等
$p=\langle\wp^1\rangle\in\Hom^0_{H(\mathcal A)}(Y,Y)$；按假设 $p$ 分裂。把
$H(\Upsilon)$ 作用于相应的图 (52)，便知对象 $Z$ 是抽象像 $\mathcal Z$ 的合法
准代表（所需的同构 $\langle t\rangle\in\Hom_{H(\mathcal Q)}(\Upsilon(Z),\mathcal Z)$
正是推论 \ref{cor:4.9} 中构造的那个）。因此 $\mathcal A$ 分裂封闭。
% [Tutor Note: 【易错】这里有两类“$Z$”：不带花体的 $Z\in\operatorname{obj}(\mathcal A)$
% 是 $p$ 的分裂对象，带花体的 $\mathcal Z\in\operatorname{obj}(\mathcal Q)$ 是抽象像模。
% “$\mathcal Z$ 被 $Z$ 准表示”指 $\Upsilon(Z)$ 与 $\mathcal Z$ 在 $H^0(\mathcal Q)$ 中同构。
% 原文用同一字母排版不同字体，抄写时务必分清。]

反过来，设 $\mathcal A$ 分裂封闭，考虑幂等元
$p=\langle\wp^1\rangle\in\Hom_{H(\mathcal A)}(Y,Y)$（再次用到引理 \ref{lem:4.6}）。
按假设，存在某个 $X\in\operatorname{obj}(\mathcal A)$，其 Yoneda 模
$\mathcal X\in\operatorname{obj}(\mathcal Q)$ 准表示 $\wp$ 的抽象像 $\mathcal Z$。由引理
\ref{lem:4.8}，
\[
  \Hom^0_{H(\mathcal Q)}(\mathcal Z,\mathcal X)
  \cong H^0(\mathcal X)(Y)\circ p
  =\Hom^0_{H(\mathcal A)}(Y,X)\circ p,
\]
故可找到 $H^0(\mathcal A)$ 中的同构 $r=\tilde r\circ p:Y\to X$。取
$Z:=\operatorname{im}(r)$ 与 $k:=p\circ r^{-1}|_Z$，便有 $k\circ r=p$ 以及
$r\circ k=\tilde r\circ p\circ r^{-1}|_Z=\mathrm{id}_Z$。因此任一 $p$ 都经由这样的
三元组 $(Z,r,k)$ 分裂，这使 $H^0(\mathcal A)$ 分裂封闭。
\end{proof}

\begin{Cor}\label{cor:4.12}
设 $\mathcal A$ 与 $\mathcal B$ 为 $A_\infty$-范畴，$\mathcal F:\mathcal A\to\mathcal B$
为 $A_\infty$-函子。则 $(\mathcal B,\mathcal F)$ 是 $\mathcal A$ 的分裂封闭，当且仅当
$(H^0(\mathcal B),H^0(\mathcal F))$ 是 $H^0(\mathcal A)$ 的分裂封闭。
\end{Cor}

由于 $H^0(\mathcal Q)$ 分裂封闭，立即可推出 $\mathcal Q$ 是分裂封闭的
$A_\infty$-范畴。$\mathcal A$ 的唯一分裂封闭便由此得出：

\begin{Pro}\label{pro:4.13}
每个 $A_\infty$-范畴 $\mathcal A$ 都有分裂封闭 $(\widetilde{\mathcal B},\widetilde{\mathcal F})$。
若另有分裂封闭 $(\mathcal B,\mathcal F)$，则存在拟等价
$\mathcal G:\mathcal B\to\widetilde{\mathcal B}$，使得
$\mathcal G\circ\mathcal F\cong\widetilde{\mathcal F}
\in\operatorname{obj}\bigl(H^0(\mathrm{fun}(\mathcal A,\widetilde{\mathcal B}))\bigr)$。
\end{Pro}

\begin{proof}
设 $\Pi\mathcal A\subset\mathcal Q$ 为对象集
\[
  \operatorname{obj}(\Pi\mathcal A):=\bigl\{\mathcal M\in\operatorname{obj}(\mathcal Q)
  \bigm|\exists\wp\ \text{使}\ \mathcal Z\cong\mathcal M\ \text{于}\ H^0(\mathcal Q)\bigr\}
\]
的全 $A_\infty$-子范畴，其中 $\mathcal Z$ 跑遍同伦幂等的抽象像。注意每个 Yoneda 模
$\mathcal Y=\Upsilon(Y)\in\operatorname{obj}(\mathcal Q)$ 都落在 $\Pi\mathcal A$ 中：
由引理 \ref{lem:4.6}，对幂等自同态 $\mathrm{id}_Y$ 存在某个 $\wp$，使得
$\mathrm{id}_Y=\langle\wp^1\rangle\in\Hom^0_{H(\mathcal A)}(Y,Y)$，而其抽象像
$\mathcal Z$ 满足
\[
  \Hom^0_{H(\mathcal A)}(Y,Y)=H^0(\mathcal Y)(Y)\circ\mathrm{id}_Y
  \cong\Hom^0_{H(\mathcal Q)}(\mathcal Z,\mathcal Y)
\]
（由引理 \ref{lem:4.8}）；在此同构下，$\mathrm{id}_Y$ 必然被送到 $H^0(\mathcal Q)$ 中的
同构 $\mathcal Z\xrightarrow{\ \sim\ }\mathcal Y$，所以 $\mathcal Y\in
\operatorname{obj}(\Pi\mathcal A)$，如所声称。

由推论 \ref{cor:4.12}，$(\widetilde{\mathcal B},\widetilde{\mathcal F})
:=(\Pi\mathcal A,\Upsilon)$ 是 $\mathcal A$ 的合法分裂封闭，因为
$(H^0(\Pi\mathcal A),H^0(\Upsilon))$ 是 $H^0(\mathcal A)$ 的分裂封闭：$H^0(\Pi\mathcal A)$
中任一幂等自同态都分裂（论证同引理 \ref{lem:4.11}），且 $H^0(\Upsilon)$ 是全忠实嵌入，
使得任一 $\mathcal M\cong\mathcal Z\in\operatorname{obj}(H^0(\Pi\mathcal A))$ 都是形如
$\langle\Upsilon^1(\wp^1)\rangle$ 的幂等自同态的像（如图 (53) 所示）。

现在考虑 $\mathcal A$ 的另一个分裂封闭 $(\mathcal B,\mathcal F)$。按假设，$\mathcal F$ 是
c-幺的，从而有（严格幺的）拉回函子 $\mathcal F^*:\mathrm{mod}(\mathcal B)\to\mathcal Q$，
如定义 \ref{def:2.37} 所述。令
$\mathcal G:=\mathcal F^*\circ\Upsilon_{\mathcal B}:\mathcal B\to\Pi\mathcal A
\subset\mathcal Q$。考察注 \ref{rem:2.38} 中那个在同构意义下交换的图，可推出
$\mathcal G\circ\mathcal F\cong\Upsilon_{\mathcal A}$ 于
$H^0(\mathrm{fun}(\mathcal A,\Pi\mathcal A))$ 中。这又说明：限制在 $\mathcal F$ 的像对象上
$\mathcal G$ 上同调全忠实（因为 $\Upsilon_{\mathcal A}$ 如此），从而由分裂封闭的定义，
它在整个 $\mathcal B$ 上也上同调全忠实。按上面同样的论证，可证明 $\Pi\mathcal A$ 的任一
对象都同构于 $\mathcal G$ 的某个像对象，这给出本质满性。于是 $\mathcal G$ 是拟等价，
由定理 2.19 可逆，从而任意两个分裂封闭只差一个拟等价。
% [Tutor Note: 【勘误】原文此处引“Theorem 2.19”，但本书并无定理 2.19；所述“拟等价可逆”
% 的结果实为定理 2.31（第 2.8 节，见 ch02）：$\mathcal G$ 是拟等价而非拟同构，故不能
% 用引理 2.19。译文照印 2.19。]
\end{proof}

\begin{Lem}\label{lem:4.14}
三角 $A_\infty$-范畴 $\mathcal A$ 的分裂封闭是三角的。
\end{Lem}

\begin{proof}
在第 \ref{sec:3.7} 节中我们论证了 $\mathcal Q$ 是三角 $A_\infty$-范畴，从而
$H^0(\mathcal Q)$ 作为普通范畴也是三角的（由命题 \ref{pro:3.47}），并且由推论
\ref{cor:4.9} 分裂封闭。在注 \ref{rem:3.55} 中我们把 $\mathcal A$ 与
$\widetilde{\mathcal Q}\subset\mathcal Q$ 等同，后者是由所有 Yoneda 模
$\mathcal Y=\Upsilon(Y)$ 生成的三角全 $A_\infty$-子范畴，其对象从而也是三角全子范畴
$H^0(\mathcal A)\subset H^0(\mathcal Q)$ 的对象。此外按定义，
$H^0(\Pi\mathcal A)\subset H^0(\mathcal Q)$ 是所有同构于某个
$\mathcal Y\in\operatorname{obj}(H^0(\mathcal A))$ 的直和项的对象构成的全子范畴
（由引理 \ref{lem:4.8}，任一抽象像的上同调模 $H(\mathcal Z)$ 都出现在某个
$H(\mathcal Y)$ 的分解中）。应用引理 \ref{lem:4.4} 可知 $H^0(\Pi\mathcal A)$ 也是三角
范畴。

但由此容易推出 $\Pi\mathcal A$ 本身是三角 $A_\infty$-范畴：它当然非空；任一对象在
$H^0(\Pi\mathcal A)$ 中都同构于某个平移对象（因为它位于 $H^0(\mathcal Q)$ 中，而后者以
$H^0(\Pi\mathcal A)$ 为全子范畴）；$H^0(\Pi\mathcal A)$ 中任一态射都可补全为正合
三角形（在 $H^0(\mathcal Q)$ 中可以，而其所有组成部分必属于三角范畴
$H^0(\Pi\mathcal A)$；于是该正合三角形在 $\Pi\mathcal A\hookrightarrow\mathcal Q$
下的原像在 $H^0(\Pi\mathcal A)$ 中仍正合，由推论 \ref{cor:3.44}）。
\end{proof}

\begin{Def}\label{def:4.15}
设 $\mathcal B$ 为分裂封闭的三角 $A_\infty$-范畴，$\mathcal A\subset\mathcal B$ 为非空的
全 $A_\infty$-子范畴。包含 $\mathcal A$、对同构封闭且分裂封闭的最小三角全
$A_\infty$-子范畴 $\widetilde{\mathcal B}\subset\mathcal B$，称为 $\mathcal B$ 的由
$\mathcal A$ \emph{分裂生成}（split-generated）的（三角）子范畴。

具体而论，$\operatorname{obj}(\widetilde{\mathcal B})$ 可以从
$\operatorname{obj}(\mathcal A)$ 出发构造：取所有可能的映射锥与平移，反复进行，最后取
所有可能的幂等自同态的像。

若 $\widetilde{\mathcal B}=\mathcal B$，则称 $\mathcal A$ \emph{分裂生成}
$\mathcal B$（此时 $\mathcal B$ 是 $\mathcal A$ 的分裂封闭；而且 $\mathcal A$ 也按
定义 \ref{def:3.22} 的意义\emph{生成} $\mathcal B$）。
\end{Def}

\begin{Def}\label{def:4.16}
设 $\mathcal A$ 为 $A_\infty$-范畴，考虑 $\mathcal A$ 的标准三角包络
$(\mathsf{Tw}\mathcal A,\iota)$ 的分裂封闭
$(\Pi(\mathsf{Tw}\mathcal A),\Upsilon_{\mathsf{Tw}\mathcal A})$（按构造，
$\Pi(\mathsf{Tw}\mathcal A)$ 由 $\mathcal A$ 分裂生成）。我们称
\[
  \mathsf D^\pi(\mathcal A):=H^0\bigl(\Pi(\mathsf{Tw}\mathcal A)\bigr)
\]
为 $\mathcal A$ 的\emph{分裂封闭导出范畴}（split-closed derived category）。
\end{Def}
% [Tutor Note: 【易错】$\mathsf D^\pi$ 的上标 $\pi$ 与基本群 $\pi_1$ 无关，它来自
% “分裂封闭”所取的幂等完备化（幂等投影）；书写时不要把它误作
% $D^{\pi_1}$。另一个常见混淆是 $\mathsf D^\pi(\mathcal A)$（分裂封闭）与
% $\mathsf D^{\mathrm b}(\mathcal A)=H^0(\mathsf{Tw}\mathcal A)$（有界导出范畴）：
% 后者不补幂等分裂。]

\begin{Cor}\label{cor:4.17}
设 $\mathcal B$ 为分裂封闭的三角 $A_\infty$-范畴，且由全 $A_\infty$-子范畴
$\mathcal A\subset\mathcal B$ 分裂生成。则存在拟等价
$\Pi(\mathsf{Tw}\mathcal A)\to\mathcal B$，从而 $H^0(\mathcal B)$ 与
$\mathsf D^\pi(\mathcal A)$ 作为分裂封闭三角范畴等价。
\end{Cor}

\begin{proof}
连同包含 $\mathcal A\hookrightarrow\mathcal B$，$\mathcal B$ 便是 $\mathcal A$ 的一个
合法三角包络兼分裂封闭。相应的唯一性结论（命题 \ref{pro:3.50} 与 \ref{pro:4.13}）
给出所需结果。
\end{proof}

由推论 \ref{cor:3.52} 直接可得：

\begin{Cor}\label{cor:4.18}
设 $\mathcal A$ 与 $\mathcal B$ 为拟等价的 c-幺 $A_\infty$-范畴。则
$\mathsf D^\pi(\mathcal A)$ 与 $\mathsf D^\pi(\mathcal B)$ 是等价的分裂封闭三角范畴。
\end{Cor}

结合最后两条推论，得到：

\begin{Cor}\label{cor:4.19}
设 $\mathcal A$ 与 $\mathcal B$ 为 c-幺 $A_\infty$-范畴，分别由全 $A_\infty$-子范畴
$\widetilde{\mathcal A}\subset\mathcal A$ 与 $\widetilde{\mathcal B}\subset\mathcal B$
分裂生成，且 $\widetilde{\mathcal A}$ 与 $\widetilde{\mathcal B}$ 拟等价。则
$\mathsf D^\pi(\mathcal A)$ 与 $\mathsf D^\pi(\mathcal B)$ 是等价的分裂封闭三角范畴。
\end{Cor}

\subsection{扭转}\label{sec:4.4}

为总结关于三角化的讨论，并给分裂生成的理论多一点立足点——我们最终要在 Fukaya 范畴的
语境中用到它——我们退后一步，抽象地讨论 $A_\infty$-模的映射锥。不过，我们选择较简略的
叙述。关于扭转的更多内容可参看 \cite{Sei08} 第 5 章。

\begin{Def}\label{def:4.20}
设 $\mathcal A$ 为 $A_\infty$-范畴，$Y\in\operatorname{obj}(\mathcal A)$。定义沿 $Y$
的\emph{抽象扭转函子}（abstract twist functor）
$\mathcal T_Y:\mathcal Q\to\mathcal Q$ 为如下 $A_\infty$-函子：
\begin{itemize}
  \item 把 $\mathcal M\in\operatorname{obj}(\mathcal Q)$ 送到
        $\mathcal C=\mathcal T_Y(\mathcal M)\in\operatorname{obj}(\mathcal Q)$，后者把
        $X\in\operatorname{obj}(\mathcal A)$ 送到
        \[
          \mathcal C(X):=\bigl(\mathcal M(Y)\otimes\hom_{\mathcal A}(X,Y)\bigr)[1]
          \oplus\mathcal M(X)\in\operatorname{obj}(\mathrm{GradVect}),
        \]
        其分次映射为
        \[
          \begin{aligned}
            \mu^1_{\mathcal C}(c_0\otimes b,c_1)
            :={}&\bigl((-1)^{|b|-1}\mu^1_{\mathcal M}(c_0)\otimes b
              +c_0\otimes\mu^1_{\mathcal A}(b),\ 
              \mu^1_{\mathcal M}(c_1)+\mu^2_{\mathcal M}(c_0,b)\bigr),\\
            \mu^d_{\mathcal C}\bigl((c_0\otimes b,c_1),a_{d-1},\dots\bigr)
            :={}&\bigl(c_0\otimes\mu^d_{\mathcal A}(b,a_{d-1},\dots),\ 
              \mu^d_{\mathcal M}(c_1,a_{d-1},\dots)\\
              &\quad+\mu^{d+1}_{\mathcal M}(c_0,b,a_{d-1},\dots)\bigr),
          \end{aligned}
        \]
        对 $d>1$ 成立。
  \item 把 $t\in\hom_{\mathcal Q}(\mathcal M_0,\mathcal M_1)$ 送到
        $\tilde t=\mathcal T^1_Y(t)\in\hom_{\mathcal Q}\bigl(\mathcal T_Y(\mathcal M_0),
        \mathcal T_Y(\mathcal M_1)\bigr)$，后者由
        \[
          \begin{aligned}
            \tilde t^1(c_0\otimes b,c_1)
            &:=\bigl((-1)^{|b|-1}t^1(c_0)\otimes b,\ t^1(c_1)+t^2(c_0,b)\bigr),\\
            \tilde t^d\bigl((c_0\otimes b,c_1),a_{d-1},\dots\bigr)
            &:=\bigl(0,\ t^d(c_1,a_{d-1},\dots)+t^{d+1}(c_0,b,a_{d-1},\dots)\bigr),
          \end{aligned}
        \]
        对 $d>1$ 成立；而当 $d>1$ 时 $\mathcal T^d_Y:=0$。
\end{itemize}
给定两个 $Y_0,Y_1\in\operatorname{obj}(\mathcal A)$，我们用
$T_{Y_0}(Y_1)\in\operatorname{obj}(\mathcal A)$ 表示
$\mathcal T_{Y_0}(\mathcal Y_1)$ 的任一准代表，并称之为 $Y_1$ 沿 $Y_0$ 的
\emph{扭转}（twist）。
\end{Def}
% [Tutor Note: 【易错】扭转模的底层分次空间是
% $(\mathcal M(Y)\otimes\hom_{\mathcal A}(X,Y))[1]\oplus\mathcal M(X)$，第一个直和项被
% 平移了一次，这正是 $\mu^1_{\mathcal C}$ 第一分量中 Koszul 符号 $(-1)^{|b|-1}$ 的来源
% （把 $b$ 从 $\hom$ 中搬到 $\mathcal M(Y)$ 前面要付 $|b|$ 次，整体再乘 $-1$）。
% 与定义 3.53 的抽象映射锥比较时，注意那里第一项是 $\hom_{\mathcal A}(X,Y_0)[1]$，
% 而这里多了一个 $\mathcal M(Y)$ 因子。]

\begin{Def}\label{def:4.21}
设 $\mathcal M,\mathcal Y=\Upsilon(Y)\in\operatorname{obj}(\mathcal Q)$，其中
$Y\in\operatorname{obj}(\mathcal A)$。
% [Tutor Note: 【勘误】原文此处印作“for $Y\in\operatorname{obj}(\mathcal Q)$”，
% 但 $\Upsilon$ 的定义域是 $\mathcal A$，故第二个 $\operatorname{obj}(\mathcal Q)$
% 应为 $\operatorname{obj}(\mathcal A)$；否则 $\mathcal Y=\Upsilon(Y)$ 无从谈起。
% 译文按 $\operatorname{obj}(\mathcal A)$ 处理，并在报告登记。]
定义\emph{抽象求值态射}（abstract evaluation
morphism）$\epsilon\in\hom^0_{\mathcal Q}\bigl(\mathcal M(Y)\otimes\mathcal Y,\mathcal M\bigr)$
为如下模同态：
\[
  \epsilon^d(c_0\otimes b,a_{d-1},\dots,a_1)
  :=\mu^{d+1}_{\mathcal M}(c_0,b,a_{d-1},\dots,a_1).
\]
则 $\mathcal T_Y(\mathcal M)=\mathrm{Cone}(\epsilon)$，如定义 \ref{def:3.53} 所示，且
它落入 $H^0(\mathcal Q)$ 中的正合三角形
\[
  \begin{tikzcd}[row sep=large, column sep=large]
    \mathcal M(Y)\otimes\mathcal Y \arrow[rr, "{\langle\epsilon\rangle}"]
      & & \mathcal M \arrow[dl, "{\langle\iota\rangle}"] \\
    & \mathcal T_Y\mathcal M \arrow[ul, "{\langle\pi\rangle{[}1{]}}"'] &
  \end{tikzcd}
\]
假设 $H(\mathcal A)$ 对有限直和与平移封闭，并给定
$Y_0,Y_1\in\operatorname{obj}(\mathcal A)$ 使
$\Hom_{H(\mathcal A)}(Y_0,Y_1)$ 有限维（从而如第 \ref{sec:3.2} 节那样可以用它作张量），
则有显然的求值态射
$\langle\mathrm{ev}\rangle:\Hom_{H(\mathcal A)}(Y_0,Y_1)\otimes\mathcal Y_0
\to\mathcal Y_1$，使得
$H(\Upsilon)\langle\mathrm{ev}\rangle=\langle\epsilon\rangle$。若扭转
$T_{Y_0}(Y_1)$ 存在，则三角形
\[
  \begin{tikzcd}[row sep=large, column sep=large]
    \Hom_{H(\mathcal A)}(Y_0,Y_1)\otimes Y_0 \arrow[rr, "{\langle\mathrm{ev}\rangle}"]
      & & Y_1 \arrow[dl, "{\langle i\rangle}"] \\
    & T_{Y_0}(Y_1) \arrow[ul, "{\langle p\rangle{[}1{]}}"'] &
  \end{tikzcd}
\]
在 $H(\Upsilon)$ 下（取 $\mathcal M=\mathcal Y_1$）映入前一三角形，故也正合（因为正合
三角形的原像仍正合；见推论 \ref{cor:3.44}）。
% [Tutor Note: 【勘误】第二个三角形中，原文把张量因子与靶对象印作对象 $Y_0,Y_1$，
% 但按上下文它们应是 Yoneda 模 $\mathcal Y_0,\mathcal Y_1$：正文紧接着说“取
% $\mathcal M=\mathcal Y_1$”，且第一分量 $\Hom_{H(\mathcal A)}(Y_0,Y_1)$ 才是纯对象性的。
% 译文照印原文的 $Y_0,Y_1$，读者按 $\mathcal Y_0,\mathcal Y_1$ 理解即可。]
\end{Def}

\begin{Rem}\label{rem:4.22}
\begin{itemize}
  \item 若 $\mathcal A$ 是三角的且 $H(\mathcal A)$ 的 Hom 空间有限维，则
        $\langle\mathrm{ev}\rangle$ 的标准映射锥为
        $\mathbb S\bigl(H(\mathcal Y_1)(Y_0)\otimes\mathcal Y_0\bigr)\oplus\mathcal Y_1$。
        以 $H(\Upsilon)$ 为竖直箭头应用引理 \ref{lem:3.3}，此锥映入
        $\mathcal T_{Y_0}(\mathcal Y_1)$，而后者总是存在。因此
        \[
          T_{Y_0}(Y_1):=\mathbb S\bigl(H(\mathcal Y_1(Y_0))\otimes\mathcal Y_0\bigr)
          \oplus\mathcal Y_1
        \]
        对所有 $Y_0,Y_1\in\operatorname{obj}(\mathcal A)$ 都有定义，从而可以把扭转函子
        $T_Y:\mathcal A\to\mathcal A$ 确定为使下图
        \[
          \begin{tikzcd}[row sep=large, column sep=large]
            \mathcal A \arrow[r, "T_Y"] \arrow[d, "\Upsilon"']
              & \mathcal A \arrow[d, "\Upsilon"] \\
            \mathcal Q \arrow[r, "\mathcal T_Y"']
              & \mathcal Q
          \end{tikzcd}
        \]
        在 $H^0(\mathrm{fun}(\mathcal A,\mathcal Q))$ 中（在同构意义下）交换的
        $A_\infty$-函子。
  \item 回忆注 \ref{rem:3.55} 中的映射
        $\widetilde\Upsilon=\iota^*\circ\Upsilon_{\mathsf{Tw}\mathcal A}:
        \mathsf{Tw}\mathcal A\to\mathcal Q$。上一个图告诉我们——在作替换
        $\mathcal A\leftarrow\mathsf{Tw}\mathcal A$ 与
        $\Upsilon\leftarrow\widetilde\Upsilon$ 之后——
        \[
          \widetilde\Upsilon\bigl(T_{Y_0}(Y_1)\bigr)
          \cong\mathcal T_{Y_0}\bigl(\widetilde\Upsilon(Y_1)\bigr)
          \in\operatorname{obj}\bigl(H^0(\mathcal Q)\bigr),
        \]
        对任意 $Y_0\in\operatorname{obj}(\mathcal A)$ 与
        $Y_1\in\operatorname{obj}(\mathsf{Tw}\mathcal A)$ 成立。
  \item 设 $\mathcal F:\mathcal A\to\mathcal B$ 为上同调全忠实的 $A_\infty$-函子，
        $Y_0,Y_1\in\operatorname{obj}(\mathcal A)$ 使得 $T_{Y_0}(Y_1)$ 与
        $T_{\mathcal F(Y_0)}(\mathcal F(Y_1))$ 都存在。则存在 $H^0(\mathcal B)$ 中的
        典范同构
        \[
          \mathcal F\bigl(T_{Y_0}(Y_1)\bigr)
          \cong T_{\mathcal F(Y_0)}\bigl(\mathcal F(Y_1)\bigr),
        \]
        而且它可延拓为
        $\mathcal F\circ T_{Y_0}\cong T_{\mathcal F(Y_0)}\circ\mathcal F
        \in\operatorname{obj}\bigl(H^0(\mathrm{fun}(\mathcal A,\mathcal B))\bigr)$。
\end{itemize}
\end{Rem}
% [Tutor Note: 【易错】第一条中映射锥的写法前后不一：先作
% $H(\mathcal Y_1)(Y_0)$，后作 $H(\mathcal Y_1(Y_0))$。二者按定义为同一个分次向量空间
% （$H$ 与“在 $Y_0$ 处取值”可交换），此处是记号混用，并非两种不同对象；抄写时取其一
% 即可，译文照印两处原样。]

下面给出一个判别法，用来判断 $\mathcal A$ 中的对象是否准表示 $\mathcal Q$ 中给定的抽象
扭转模。

\begin{Lem}\label{lem:4.23}
设 $\mathcal A$ 为 $A_\infty$-范畴，$Y_0,Y_1,Y_2\in\operatorname{obj}(\mathcal A)$。
假设存在上闭链 $c\in\hom_{\mathcal A}(Y_1,Y_2)$ 以及映射
$k:\hom_{\mathcal A}(Y_0,Y_1)\to\hom_{\mathcal A}(Y_0,Y_2)[-1]$，使得对一切
$c_0\in\hom_{\mathcal A}(Y_0,Y_1)$ 有
\[
  \mu^1_{\mathcal A}(k(c_0))-k\bigl(\mu^1_{\mathcal A}(c_0)\bigr)
  +\mu^2_{\mathcal A}(c,c_0)=0.
\]
则如下模同态 $t:\mathcal T_{Y_0}(\mathcal Y_1)\to\mathcal Y_2$ 良定义：
\[
  \begin{aligned}
    t^d\bigl((c_0\otimes b,c_1),a_{d-1},\dots,a_1\bigr)
    :={}&\mu^{d+1}_{\mathcal A}(c,c_1,a_{d-1},\dots)
      +\mu^{d+1}_{\mathcal A}(k(c_0),b,a_{d-1},\dots)\\
      &+\mu^{d+2}_{\mathcal A}(c,c_0,b,a_{d-1},\dots),
  \end{aligned}
\]
它满足 $\langle\mu^2_{\mathcal Q}(t,\iota)\rangle=\langle\Upsilon^1(c)\rangle
\in\Hom_{H(\mathcal Q)}(\mathcal Y_1,\mathcal Y_2)$；并且——在给定适当的附加数据时
（见 \cite{Sei08} 引理 5.3）——它在 $H^0(\mathcal Q)$ 中诱导同构
$\mathcal T_{Y_0}(\mathcal Y_1)\cong\mathcal Y_2$。
\end{Lem}
% [Tutor Note: 【易错】条件中的三项缺一不可：$\mu^1_{\mathcal A}(k(c_0))$ 是 $k(c_0)$
% 的微分，$-k(\mu^1_{\mathcal A}(c_0))$ 是 $k$ 与 $c_0$ 的微分交换失败项，而
% $\mu^2_{\mathcal A}(c,c_0)$ 记录 $c$ 与 $c_0$ 的乘法；把最后一项误写为
% $\mu^2_{\mathcal A}(c_0,c)$ 会破坏 $t$ 的上闭链性。]

当然，扭转运算可以迭代使用。考虑 Hom 空间全都有限维的三角 $A_\infty$-范畴
$\mathcal A$。取集合 $\{Y_1,\dots,Y_m\}\subset\operatorname{obj}(\mathcal A)$，并从
$X:=X_m\in\operatorname{obj}(\mathcal A)$ 出发，依次定义
$X_{k-1}:=T_{Y_k}(X_k)\in\operatorname{obj}(\mathcal A)$，它们各自落入形如定义
\ref{def:4.21} 的正合三角形。利用三角范畴的公理 (T3)，可以把这些三角形复合起来，
得到另一个正合三角形
\[
  \begin{tikzcd}[row sep=large, column sep=large]
    Y \arrow[rr] & & X \arrow[dl, "{\langle i_1\rangle\circ\dots\circ\langle i_m\rangle}"] \\
    & (T_{Y_1}\circ\dots\circ T_{Y_m})(X) \arrow[ul, "{[}1{]}"'] &
  \end{tikzcd}
\]
其中 $Y\in\mathcal A$ 由诸空间 $\Hom_{H(\mathcal A)}(Y_k,X_k)\otimes Y_k$ 适当构造，
而 $\langle i_k\rangle:X_k\to X_{k-1}$ 是典范包含。正如注 \ref{rem:3.48} 中所见，
$\langle i_1\rangle\circ\dots\circ\langle i_m\rangle$ 的消没蕴含
$X\oplus(T_{Y_1}\circ\dots\circ T_{Y_m})(X)[-1]\cong Y$。以此为基础可以证明如下结果。

\begin{Cor}\label{cor:4.24}
设 $\mathcal A$ 为 Hom 空间全都有限维的三角 $A_\infty$-范畴。
\begin{enumerate}[label=\textup{\roman*}.,leftmargin=*]
  \item 假设存在 $Y_1,\dots,Y_m\in\operatorname{obj}(\mathcal A)$，使得对一切
        $X\in\operatorname{obj}(\mathcal A)$ 有
        $(T_{Y_1}\circ\dots\circ T_{Y_m})(X)\cong 0
        \in\operatorname{obj}(H^0(\mathcal A))$（即 $H^0(\mathcal A)$ 的零对象）。
        则 $\{Y_1,\dots,Y_m\}$ 生成 $\mathcal A$，而 $\mathcal A$ 其实是分裂封闭的。
  \item 假设 $\mathcal A$ 已分裂封闭，且存在
        $Y_1,\dots,Y_m\in\operatorname{obj}(\mathcal A)$ 与 $\sigma\neq0$，使得对一切
        $X\in\operatorname{obj}(\mathcal A)$ 有
        $(T_{Y_1}\circ\dots\circ T_{Y_m})(X)\cong S^\sigma X
        \in\operatorname{obj}(H^0(\mathcal A))$。则 $\{Y_1,\dots,Y_m\}$ 分裂生成
        $\mathcal A$。
\end{enumerate}
\end{Cor}
% [Tutor Note: 【含糊】原文 (ii) 的假设写作“$\mathcal A$ is readily split-closed”，
% “readily”在此语义不明；从证明用法看应读作“已经分裂封闭（already split-closed）”。
% 译文按“已分裂封闭”处理，并在此说明原文措辞。若按字面“轻易地”理解，则该假设无法
% 与推论 (i) 中的“$\mathcal A$ 其实是分裂封闭的”相区分。]

\begin{Rem}\label{rem:4.25}
在考虑 $\mathsf{Tw}\mathcal A$ 中的扭复形时（其中 $\mathcal A$ 取为严格幺的），我们其实
可以改进对扭转函子的描述。

取 $Y_0,Y_1\in\operatorname{obj}(\mathsf{Tw}\mathcal A)$，考虑
$\bigl(\hom_{\mathsf{Tw}\mathcal A}(Y_0,Y_1),\partial:=\mu^1_{\mathsf{Tw}\mathcal A}\bigr)$，
其对偶复形为
$\hom_{\mathsf{Tw}\mathcal A}(Y_0,Y_1)^\vee:=
\hom_{\mathsf{Tw}\mathcal A}\bigl(\hom_{\mathsf{Tw}\mathcal A}(Y_0,Y_1),Y_1\bigr)$，
微分 $\partial^\vee$ 由
$(\partial^\vee(a^\vee))(b):=(-1)^{|a^\vee|+1}a^\vee(\partial(b))$ 给出，其中
$a^\vee\in\hom_{\mathsf{Tw}\mathcal A}(Y_0,Y_1)^\vee$、
$b\in\hom_{\mathsf{Tw}\mathcal A}(Y_0,Y_1)$。可以如定义 \ref{def:3.33} 那样取张量积：
\[
  \begin{aligned}
    \hom_{\mathsf{Tw}\mathcal A}\bigl(\hom_{\mathsf{Tw}\mathcal A}(Y_0,Y_1)
      \otimes Y_0,Y_1\bigr)
    &=\hom_{\mathsf{Tw}\mathcal A}(Y_0,Y_1)^\vee
      \otimes\hom_{\mathsf{Tw}\mathcal A}(Y_0,Y_1)\\
    &\cong\mathrm{End}_{\mathbb K}\bigl(\hom_{\mathsf{Tw}\mathcal A}(Y_0,Y_1)\bigr),
  \end{aligned}
\]
其微分为 $D:a^\vee\otimes b\mapsto a^\vee(\cdot)b$，而
$\mathrm{ev}\cong\mathrm{id}\in\mathrm{End}_{\mathbb K}
\bigl(\hom_{\mathsf{Tw}\mathcal A}(Y_0,Y_1)\bigr)$ 是关于它的上闭链。因此可以考察扭
映射锥 $\mathrm{Cone}(\mathrm{ev})$，它在 $H(\Upsilon)$ 下变成
$\mathrm{Cone}(\epsilon)=\mathcal T_{Y_0}(\mathcal Y_1)$（参看定义 \ref{def:4.21}）。
这促使我们置
\[
  T_{Y_0}(Y_1):=\mathrm{Cone}(\mathrm{ev})
  =\bigoplus_{i\in I}S^{-|b^i|+1}Y_1\oplus Y_0,
\]
其中 $\{b^i\}_{i\in I}$ 是 $\hom_{\mathcal A}(Y_0,Y_1)$ 的一组基。
% [Tutor Note: 【勘误】原文此处印作“basis of $\hom_{\mathcal A}(Y_0,Y_1)$”，
% 但本注开头取的是 $Y_0,Y_1\in\operatorname{obj}(\mathsf{Tw}\mathcal A)$，故基应取自
% $\hom_{\mathsf{Tw}\mathcal A}(Y_0,Y_1)$；只有回到 $\mathcal A\subset\mathsf{Tw}\mathcal A$
% 时两者才一致。译文照印 $\hom_{\mathcal A}$。]
\end{Rem}
% [Tutor Note: 【背景】这里把扭复形中的“扭转函子”具体化为一个映射锥：对偶复形
% $\hom(Y_0,Y_1)^\vee=\hom(\hom(Y_0,Y_1),Y_1)$ 把 $\hom$ 空间看作复形，
% $\mathrm{ev}$ 是求值同态；于是 $T_{Y_0}(Y_1)$ 是 $\hom_{\mathcal A}(Y_0,Y_1)$
% 的一组基对应的若干平移 $Y_1$ 与一个 $Y_0$ 的直和。这与定义 4.21 的抽象锥在
% $H(\Upsilon)$ 下一致，见 \cite{Sei08} 第 5 章。]

最后，我们提及扭转运算的伴随版本。

\begin{Def}\label{def:4.26}
设 $\mathcal A$ 为 $A_\infty$-范畴，$Y_0,Y_1\in\operatorname{obj}(\mathcal A)$。
$A_\infty$-模 $\mathcal C^\vee\in\operatorname{obj}(\mathcal Q)$ 由下式给出：
\[
  \mathcal C^\vee(X):=\hom_{\mathcal A}(X,Y_1)
  \oplus\Hom_{\mathbb K}\bigl(\hom_{\mathcal A}(Y_1,Y_0),
  \hom_{\mathcal A}(X,Y_0)\bigr)[-1],
\]
\[
  \begin{aligned}
    \mu^1_{\mathcal C^\vee}(c,\beta)
    &:=\bigl(\mu^1_{\mathcal A}(c),\
      \mu^1_{\mathcal A}\circ\beta+(-1)^{|\beta|-1}\beta\circ\mu^1_{\mathcal A}
      +\mu^2_{\mathcal A}(\cdot,c)\bigr),\\
    \mu^d_{\mathcal C^\vee}\bigl((c,\beta),a_{d-1},\dots\bigr)
    &:=\bigl(\mu^d_{\mathcal A}(c,a_{d-1},\dots),\
      \mu^d_{\mathcal A}(\beta(\cdot),a_{d-1},\dots)
      +\mu^{d+1}_{\mathcal A}(\cdot,c,a_{d-1},\dots)\bigr),
  \end{aligned}
\]
对 $d>1$、$X\in\operatorname{obj}(\mathcal A)$ 成立。
\end{Def}

我们记 $\mathcal C^\vee$ 的任一准代表为 $T^\vee_{Y_0}(Y_1)$，并称之为 $Y_1$ 沿 $Y_0$ 的
\emph{伴随扭转}（adjoint twist）。可以证明，对一切
$Y\in\operatorname{obj}(\mathcal A)$ 有 $T_{Y_0}\bigl(T^\vee_{Y_0}(Y)\bigr)\cong Y$ 于
$H^0(\mathcal A)$ 中（事实上，$T^\vee_{Y_0}$ 确实是 $T_{Y_0}$ 在
$H^0(\mathrm{fun}(\mathcal A,\mathcal A))$ 中的逆，至多差同构）。

与通常的扭转一样，取 $\mathcal A$ 为 Hom 空间有限维的三角范畴时，
$T^\vee_{Y_0}(Y_1)$ 总有定义，从而落入“对偶化”后得到的正合三角形：
\[
  \begin{tikzcd}[row sep=large, column sep=large]
    \Hom_{H(\mathcal A)}(Y_1,Y_0)^\vee\otimes Y_0
      \arrow[dr, "{\langle p^\vee\rangle{[}1{]}}"']
    & & Y_1 \arrow[ll, "{\langle\mathrm{ev}^\vee\rangle}"'] \\
    & T^\vee_{Y_0}(Y_1) \arrow[ur, "{\langle i^\vee\rangle}"'] &
  \end{tikzcd}
\]
% [Tutor Note: 【易错】此三角形与前文的正合三角形方向相反（对偶化使箭头反转）：
% $\langle\mathrm{ev}^\vee\rangle$ 指向左、$\langle p^\vee\rangle[1]$ 从左上指向下、
% $\langle i^\vee\rangle$ 从下指向右上。图中张量因子按原文印作对象 $Y_0$，理应按
% Yoneda 模 $\mathcal Y_0$ 理解（与定义 4.21 第二个三角形的记号问题同源）。]

再次取 $\mathcal A$ 为严格幺的，$T^\vee_{Y_0}(Y_1)$ 可以在 $\mathsf{Tw}\mathcal A$ 中
等同于“对偶”求值态射 $\langle\mathrm{ev}^\vee\rangle$ 的映射锥（与上一注平行）。

